\chapter{よくある質問と罠}
\label{app:faq}

この付録は1冊を学習し、読者がよく経験する問題と質問をまとめたものである。各項目は、実際の実装中に発生する可能性がある具体的な状況を扱います。

\section{環境関連}

\subsection{Q: PyTorch インストール後 CUDAが認識されません。}

\textbf{A}：2つのことを確認してください。まず、NVIDIAドライバがインストールされているか（\code{nvidia-smi}実行）。第二に、PyTorchバージョンとCUDAバージョンが互換性があるか。たとえば、CUDA 12.1用のPyTorchをCUDA 11.8システムにインストールすると認識されません。

\begin{lstlisting}[style=bashstyle]
# CUDA バージョン確認
nvidia-smi | grep "CUDA Version"

# PyTorch 再インストール（CUDA 12.1用）
pip uninstall torch
pip install torch --index-url https://download.pytorch.org/whl/cu121
\end{lstlisting}

\subsection{Q: 学習が遅すぎる}

\textbf{A}：以下を確認してください：
\begin{enumerate}
\item GPUを使用しているか（\code{device="cuda"}設定）。
\item バッチサイズが小さすぎないか（GPU使用率も低い）。
\item データのロードがボトルネックか (\code{num\_workers}増加)。
\item 混合精度（BF16）を使用しているか。
\end{enumerate}

\section{トークナイザー関連}

\subsection{Q: BPEで学んだトークナイザー ``こんにちは''奇妙なトークン化}

\textbf{A}：韓国語がコーパスに十分に含まれていないため。 BPEは頻度ベースであり、まれな文字はバイトに分解されます。解決策：
\begin{itemize}
\item 韓国語コーパスの割合を高めること。
\item 語彙サイズを増やすこと（韓国語音節約2万個＋サブワード）。
\item SentencePieceのようなプロのトークナイザーを使用。
\end{itemize}

\subsection{Q: encode/decode ラウンドトリップが失敗しました。}

\textbf{A}：空白の処理を確認してください。 BPEの「」文字がデコード時に空白に変換されない可能性があります。私たちの実装では、\code{\_postprocess}メソッドがこれを担当します。

\begin{lstlisting}[language=Python]
def _postprocess(self, text):
    text = text.replace("▁", " ")
    return text.strip()
\end{lstlisting}

\subsection{Q: 特殊トークン IDが変わるとどうなりますか？}

\textbf{A}：致命的。モデルは\code{<pad>=0, <bos>=1, <eos>=2, <unk>=3}を想定して学習されました。トークナイザーを再学習しながら特殊トークンIDが変わると、モデルの重みはもはや有効ではありません。常に特別なトークンIDを固定してください。

\section{モデル関連}

\subsection{Q: アテンション出力 NaNこれが出ます。}

\textbf{A}：3つの原因が一般的です：
\begin{enumerate}
\item\textbf{ソフトマックス飽和}：ロジットが大きすぎてsoftmaxが0/1に飽和しました。スケーリング（$\sqrt{d\_k}$）確認。
\item\textbf{学習率が大きすぎる}：グラデーション爆発。学習率を1/10に減らしてください。
\item\textbf{FP16オーバーフロー}：アクティブ値がFP16の範囲（$\pm 65504$）を超えています。 BF16に切り替えます。
\end{enumerate}

\begin{lstlisting}[language=Python]
# NaN 確認
if torch.isnan(loss):
    print("NaN detected! Check learning rate and gradients.")
    print(f"Max grad: {max(p.grad.abs().max() for p in model.parameters() if p.grad is not None)}")
\end{lstlisting}

\subsection{Q: モデルは同じ単語を繰り返します。}

\textbf{A}: サンプリング戦略を確認してください。 Greedyは繰り返しループに陥りやすいです。 Top-P（P = 0.9）に切り替えるか、Temperatureを0.7$\sim$0.9に設定します。繰り返しペナルティの追加も効果的。

\begin{lstlisting}[language=Python]
# 繰り返しペナルティ(簡素化)
def apply_repetition_penalty(logits, generated_ids, penalty=1.2):
    for token_id in set(generated_ids.tolist()):
        logits[token_id] /= penalty
    return logits
\end{lstlisting}

\subsection{Q: コンテキスト長を超えるとエラーが発生します}

\textbf{A}：モデルの\code{context\_length}を超える入力を入れると埋め込みフェーズでエラー。\code{generate}関数からスライドウィンドウに切り取ることを確認してください。

\begin{lstlisting}[language=Python]
if input_ids.shape[1] > context_length:
    x = input_ids[:, -context_length:]  # 最新のものだけを維持
\end{lstlisting}

\section{学習関連}

\subsection{Q: 損失は​​減りません。}

\textbf{A}：以下を確認してください：
\begin{enumerate}
\item 学習率が小さすぎる(1e-6以下)、大きすぎる(1e-2以上).
\item データが少なすぎます（最小数千のトークンが必要です）。
\item オプティマイザが無効です（AdamWのベータ、epsを確認）。
\item グラデーションが流れない (requires\_grad チェック)。
\end{enumerate}

\begin{lstlisting}[language=Python]
# グラデーションフローの確認
for name, p in model.named_parameters():
    if p.requires_grad and p.grad is None:
        print(f"No gradient for {name}")
    elif p.grad is not None and p.grad.abs().max() < 1e-10:
        print(f"Vanishing gradient for {name}")
\end{lstlisting}

\subsection{Q: 学習の初期に損失が発散します。}

\textbf{A}：Warmupが不足しているか、初期化が間違っている可能性があります。以下を試してください：
\begin{itemize}
\item Warmup ステップを増やす (100$\sim$1000).
\item グラデーションクリッピングチェック（\code{grad\_clip=1.0}）。
\item 初期化標準偏差を減らす（0.02$\to$0.01）。
\end{itemize}

\subsection{Q: チェックポイントをロードすると他の結果が表示されます。}

\textbf{A}：2つの原因：
\begin{enumerate}
\item\textbf{シード未設定}:\code{set\_seed(42)}を呼び出したことを確認します。
\item\textbf{ドロップアウト}：モデルを\code{eval()}モードに切り替えたことを確認します。
\item\textbf{CuDNN 非決定論}:\code{torch.backends.cudnn.deterministic = True}.
\end{enumerate}

\section{分散学習関連（2巻プレビュー）}

\subsection{Q: DDPで学習すると速度がむしろ遅くなります。}

\textbf{A}：通信ボトルネックである可能性。以下を確認してください。
\begin{itemize}
\item バッチサイズが小さすぎ、演算/通信比が低い。
\item NVLinkではなく、低速通信（PCIeなど）を使用。
\item グラデーション同期がステップごとに起こります (gradient accumulation に減少)。
\end{itemize}

\subsection{Q: FSDPから OOMこれは私です}

\textbf{A}：アクティブメモリが不足。有効チェックポイントをオンにするか、バッチサイズを減らします。 CPUオフロードも考慮。

\section{デバッグのヒント}

\subsection{段階的に検証する}

複雑なバグは小さな単位で割って検証。トークナイザーのみ単独テスト、アテンションのみ単独テストなど。

\begin{lstlisting}[language=Python]
# トルクナイザー単独テスト
tok = BPETokenizer()
tok.train(["hello world"], vocab_size=50)
print(tok.encode("hello"))
print(tok.decode(tok.encode("hello")))

# アテンション単独テスト
attn = MultiHeadAttention(config)
x = torch.randn(1, 4, 32)
out = attn(x)
print(out.shape, out.dtype, torch.isnan(out).any())
\end{lstlisting}

\subsection{shapeを常に出力する}

バグの 80\% は shape 不一致。あらゆる段階で形状を出力する習慣。

\begin{lstlisting}[language=Python]
print(f"input shape: {x.shape}")
print(f"after embedding: {emb.shape}")
print(f"after attention: {attn_out.shape}")
\end{lstlisting}

\subsection{小さな入力で始める}

バッチ1、シーケンス4で始まり、動作を確認してサイズを大きくする。大きな入力でデバッグすると出力が多すぎて混乱。

\section{パフォーマンス最適化のヒント}

\subsection{混合精度の使用}

\begin{lstlisting}[language=Python]
from torch.cuda.amp import autocast, GradScaler

scaler = GradScaler()
with autocast(dtype=torch.bfloat16):
    logits = model(inputs)
    loss = criterion(logits, targets)
scaler.scale(loss).backward()
scaler.step(optimizer)
scaler.update()
\end{lstlisting}

BF16はFP16とは異なり、スケーラを必要としません。 A100以上で推奨。

\subsection{torch.compile}

PyTorch 2.0の\code{torch.compile}でモデルをコンパイルすると、1.5$\sim$の2倍の速度が向上します。

\begin{lstlisting}[language=Python]
model = GPT(config)
model = torch.compile(model)  # コンパイル
\end{lstlisting}

\subsection{プロファイリング}

\begin{lstlisting}[language=Python]
with torch.profiler.profile(activities=[torch.profiler.ProfilerActivity.CPU,
                                         torch.profiler.ProfilerActivity.CUDA]) as prof:
    # 学習コード
prof.export_chrome_trace("trace.json")
\end{lstlisting}

Chromeブラウザで\code{chrome://tracing}で開き、ボトルネック分析。
